\section{模块 1 · 任务与起点：上一步交给我们什么，这一步要做出什么}

\subsection{我们手里的数据长什么样}

整个项目的原始数据是\textbf{机载激光雷达（LiDAR）点云}。它的获取方式一句话就能说清：
飞机沿着测区上空飞行，向地面连续发射激光脉冲，测量每束激光打到物体反射回来的时间，
从而算出打中位置的三维坐标。每束激光的落点就是一个``点''，
千万个点合在一起就是``点云''——可以想象成\textbf{给整个城市拍了一张由无数小点组成的三维照片}：
屋顶上有点、地面上有点、树冠上也有点。

每个点至少带三个数：$x$（东西向位置）、$y$（南北向位置）、$z$（高度——注意它是
相对统一高度基准面测的，测绘上叫``高程''，不是``离地面多高''）。
本项目数据的坐标单位是\textbf{美制测量英尺}（ftUS，1 ftUS $\approx$ 0.3048 米，
这是美国政府测绘数据的惯例单位）；我们在计算时统一换算成米。

\begin{numbox}[：本项目的原始数据]
研究区为哥伦布市近东区约 2.76 平方公里的住宅区；
数据来自 USGS（美国地质调查局）2019 年航飞，共 19 块瓦片（tile，即数据分块文件），
合计 \textbf{76{,}351{,}005 个点}，平均密度约 \textbf{30 点/平方米}
——如果把点想象成均匀铺开，相邻点相距约 18 厘米
（实际分布并不均匀，这只是个数量级的直觉）。
\end{numbox}

\subsection{上一步（RandLA-Net）已经做了什么：语义分割}

原始点云只有坐标，没有``身份''：计算机不知道哪个点是屋顶、哪个点是树。
我们的整套处理流程（下文简称\textbf{管线}，pipeline，一步接一步的加工流水线）第 1 步
用一个叫 RandLA-Net 的\textbf{深度学习网络}（一类从人工标注的示例中自动学出判断规律的程序）
做了\textbf{语义分割}（semantic segmentation）——直白地说，就是\textbf{给每一个点贴一张类别标签}。
本项目用四类标签（编号沿用美国摄影测量与遥感学会 ASPRS 制定的行业标准类别码）：

\begin{center}
\begin{tabular}{cll}
\toprule
类别码 & 含义 & 占比 \\
\midrule
2 & 地面（马路、草坪、裸土） & —— \\
5 & 植被（树冠、灌木） & —— \\
6 & \textbf{建筑（屋顶、部分墙面）} & 约 863 万点，占全部点的 11.3\% \\
1 & 其他（车辆、电线等零散物） & —— \\
\bottomrule
\end{tabular}
\end{center}

于是我们得到了 19 块``每个点都带类别标签''的点云。本讲义的一切都从这份结果出发：
\textbf{把标签为 6 的约 863 万个点拿出来加工}。下文把它们简称``建筑点''，
但请记住：这只是模型打的标签，里面混着被错标的桥、树冠、电线
——清理它们正是步骤 3 的任务。另外两类标签也有用：标签 2（地面）在第 3 步会用到
（用来算``离地高度''），标签 5（植被）在后面的追溯修正中用来抓混进来的树冠。

\subsection{本步的任务：从``语义''到``实例''}

语义分割只回答``是什么''，不回答``是哪一个''。全区两千多栋楼的屋顶点，标签全都是 6，
彼此之间没有任何区分。而后续的屋顶重建必须\textbf{一栋一栋地做}
（每栋楼单独提取屋顶平面、单独拼模型），所以必须先把建筑点\textbf{按栋分开}
——这一步叫\textbf{实例分割}（instance segmentation），行业里也叫\textbf{建筑单体化}。
``实例''就是``具体的一个个体''：同样是``建筑''这个类别，1 号楼和 2 号楼是两个不同的实例。

\begin{figure}[H]
\centering
\begin{tikzpicture}[scale=0.82, every node/.style={font=\small}]
% ---- 左：语义分割 ----
\begin{scope}
\node[font=\bf\small] at (2.6,3.1) {语义分割的结果（第 1 步）};
\fill[blue!30] (0,0) rectangle (1.5,1.2);
\fill[blue!30] (2.0,0.1) rectangle (3.4,1.4);
\fill[blue!30] (3.9,0) rectangle (5.2,1.1);
\node at (0.75,0.6) {建筑};
\node at (2.7,0.75) {建筑};
\node at (4.55,0.55) {建筑};
\node[text width=5.2cm, align=center] at (2.6,-0.8)
  {三栋楼颜色相同：只知道``都是建筑''\\ 分不清彼此};
\end{scope}
% ---- 右：实例分割 ----
\begin{scope}[xshift=7.3cm]
\node[font=\bf\small] at (2.6,3.1) {实例分割的目标（本步）};
\fill[red!40] (0,0) rectangle (1.5,1.2);
\fill[green!45] (2.0,0.1) rectangle (3.4,1.4);
\fill[orange!55] (3.9,0) rectangle (5.2,1.1);
\node at (0.75,0.6) {1 号楼};
\node at (2.7,0.75) {2 号楼};
\node at (4.55,0.55) {3 号楼};
\node[text width=5.2cm, align=center] at (2.6,-0.8)
  {每栋楼一个编号（instance\_id）\\ 楼内所有点共享这个编号};
\end{scope}
\end{tikzpicture}
\caption{语义分割 vs.\ 实例分割：本步的任务就是从左图走到右图。}
\label{fig:sem-vs-inst}
\end{figure}

本步做完后，每个建筑点会多带一个整数编号 \texttt{instance\_id}：
同一栋楼的点编号相同，不同楼编号不同。后续所有步骤（屋顶平面提取、模型装配、精度评价）
都以这个编号为``楼''的定义。

\subsection{为什么这件事不是``顺手就能做''的}

直觉上似乎很简单：``挨在一起的点算一栋楼''。但真实数据里有两类拦路虎，
分别对应本步的两个子任务：

\textbf{拦路虎一：楼和楼``粘''在一起（对应步骤 2：单体分割）。}
本区是独栋住宅区，相邻两栋楼之间的侧院间隙常常只有 1.5--2.5 米。
点云是离散的、屋檐还会向外挑出，两栋楼的点在空间上非常接近。
而计算机做分组时，总要先定一个``多近才算连在一起''的尺度：尺度设小了，
一栋楼自己会碎成几块；设大了，隔着窄侧院的两栋楼就被当成一体。
一旦这个尺度处理不当（后面会看到 v1 版本的真实教训），整排房子会被连成``一栋''，
这种错误叫\textbf{欠分割}（under-segmentation，该分开的没分开）。
反过来，一栋楼如果被拆成好几块，叫\textbf{过分割}（over-segmentation）。
单体分割的全部功夫，就是在这两种错误之间走钢丝。

\textbf{拦路虎二：混进来的``假建筑''（对应步骤 3：噪声清除）。}
RandLA-Net 不是完美的，它会把一些``长得像建筑''的东西错标成类别 6：
高架桥面、立交匝道、成排的树冠、电力线。这些东西如果不清理，
会一路混到最后变成一个个``假房子模型''。而且越到后面越难发现
——桥面本身就是一块平整的板，后续算法用数学平面去近似它时，比近似真屋顶还顺利，
根本不会报错。所以必须\textbf{在程序开始逐栋重建之前、以``栋''为单位}把它们拦下来，
这道关卡我们称为 Gate A（gate＝闸门，A 表示它是全管线的第一道质量闸）。

\begin{plainbox}[：一个比喻串起两个子任务]
把 863 万个``被标成建筑''的点想象成一大袋混在一起的乐高零件。步骤 2 是\textbf{分拣}：
按空间位置把零件分成 2{,}296 小袋，每袋暂时当作``一栋候选建筑''；步骤 3 是\textbf{质检}：
逐袋检查，把``根本不是楼''的袋子（桥、树、电线）整袋挑出去。
最后剩下 2{,}267 袋通过质检的候选建筑，交给下一步拼屋顶。目标是``一袋恰好一栋楼''
——至于做到了几分，靠人工抽查与最终的精度评价来回答，而不是靠这个数字本身。
（2{,}296 与 2{,}267 怎么来的，模块 4--6 会逐项对账。）
\end{plainbox}

\subsection{本步在整条管线中的位置}

\begin{center}
\begin{tikzpicture}[node distance=0.55cm, every node/.style={font=\small}]
\tikzstyle{stage}=[draw, rounded corners, minimum height=0.9cm, text width=2.5cm, align=center]
\tikzstyle{cur}=[stage, fill=blue!15, draw=blue!60!black, thick]
\node[stage] (s1) {第 1 步\\ 语义分割};
\node[cur, right=of s1] (s23) {\textbf{第 2--3 步}\\ \textbf{单体化+清噪}};
\node[stage, right=of s23] (s45) {第 4--6 步\\ 屋顶平面·线索·轮廓};
\node[stage, right=of s45] (s7) {第 7 步起\\ 模型装配$\cdots$};
\draw[-{Stealth}] (s1) -- (s23);
\draw[-{Stealth}] (s23) -- (s45);
\draw[-{Stealth}] (s45) -- (s7);
\end{tikzpicture}
\end{center}

输入：19 块带语义标签的点云（\texttt{01\_semantic\_randlanet/results/preds/}）；
输出：一份带 \texttt{instance\_id} 的干净建筑点云文件 \texttt{buildings\_clean.laz}
（\texttt{.laz} 是点云的压缩文件格式；文件存放在 \texttt{02\_03\_instance\_seg/output\_v2/}），
外加一张逐实例的``审计表''（记录每个实例的属性和去留原因，模块 5--6 详述）。

\begin{plainbox}[：模块 1 小结]
\begin{itemize}
\item 起点：约 7{,}635 万点的点云，每点已带类别标签；其中被标为建筑的约 863 万个（11.3\%）。
\item 任务：给建筑点按栋编号（实例分割），并剔除``假建筑''（Gate A 质检）。
\item 两大风险：欠分割（楼粘成一条）与假建筑混入（桥/树/电线）；
      下面两个模块先讲我们参考了谁的方法、为什么这么选。
\end{itemize}
\end{plainbox}
